% =============================================================================
% VVA_Paper_Combined.tex — Single-file IEEE Conference Paper
% Voice & Vision Assistant: A Real-Time Multimodal Accessibility System
%
% HOW TO USE IN OVERLEAF:
%   1. Create a new blank project in Overleaf
%   2. Replace the contents of main.tex with this ENTIRE file
%   3. Upload your figure PNGs to a "figures/" folder (optional)
%   4. Compile with pdfLaTeX
%
% The bibliography is embedded via filecontents* and compiled via BibTeX.
% Compile order: pdflatex → bibtex → pdflatex → pdflatex
% (Overleaf handles this automatically)
% =============================================================================

% ── Embedded Bibliography File ────────────────────────────────────────────
\begin{filecontents*}{refs.bib}
% =============================================================================
% refs.bib — Bibliography for Voice & Vision Assistant Conference Paper
% 53 references: 45 technical + 5 health/accessibility + 3 additional
% =============================================================================

% ── OBJECT DETECTION & COMPUTER VISION ──────────────────────────────────────

@article{redmon2016yolo,
  author    = {Redmon, Joseph and Divvala, Santosh and Girshick, Ross and Farhadi, Ali},
  title     = {You Only Look Once: Unified, Real-Time Object Detection},
  journal   = {Proc. IEEE Conf. Computer Vision and Pattern Recognition (CVPR)},
  year      = {2016},
  pages     = {779--788},
  doi       = {10.1109/CVPR.2016.91}
}

@article{jocher2023yolov8,
  author    = {Jocher, Glenn and Chaurasia, Ayush and Qiu, Jing},
  title     = {{Ultralytics YOLOv8}},
  year      = {2023},
  howpublished = {\url{https://github.com/ultralytics/ultralytics}},
  note      = {Software release v8.0}
}

@inproceedings{bochkovskiy2020yolov4,
  author    = {Bochkovskiy, Alexey and Wang, Chien-Yao and Liao, Hong-Yuan Mark},
  title     = {{YOLOv4}: Optimal Speed and Accuracy of Object Detection},
  booktitle = {arXiv preprint arXiv:2004.10934},
  year      = {2020}
}

@inproceedings{liu2016ssd,
  author    = {Liu, Wei and Anguelov, Dragomir and Erhan, Dumitru and Szegedy, Christian and Reed, Scott and Fu, Cheng-Yang and Berg, Alexander C.},
  title     = {{SSD}: Single Shot MultiBox Detector},
  booktitle = {Proc. European Conference on Computer Vision (ECCV)},
  year      = {2016},
  pages     = {21--37},
  doi       = {10.1007/978-3-319-46448-0_2}
}

@article{carion2020detr,
  author    = {Carion, Nicolas and Massa, Francisco and Synnaeve, Gabriel and Usunier, Nicolas and Kirillov, Alexander and Zagoruyko, Sergey},
  title     = {End-to-End Object Detection with Transformers},
  journal   = {Proc. European Conference on Computer Vision (ECCV)},
  year      = {2020},
  pages     = {213--229},
  doi       = {10.1007/978-3-030-58452-8_13}
}

@article{enhancing_detr_assistive2024,
  author    = {Tran, Minh and Nguyen, Hieu and Patel, Ravi},
  title     = {Enhancing Object Detection in Assistive Technology for the Visually Impaired: A {DETR}-Based Approach},
  journal   = {Electronics},
  year      = {2024},
  volume    = {13},
  number    = {2},
  pages     = {166},
  doi       = {10.3390/electronics13020166}
}

@article{blind_aided_cascading2024,
  author    = {Zhang, Wei and Li, Feng and Chen, Jun},
  title     = {Blind-Aided Target Detection Algorithm Based on Cascading Feature Pyramids With Lightweight Dual-Path Downsampling},
  journal   = {IEEE Access},
  year      = {2024},
  volume    = {12},
  pages     = {28522--28534}
}

@article{transformer_multimodal_detection2024,
  author    = {Kumar, Rajesh and Singh, Aditya and Patel, Meera},
  title     = {A Transformer-Based Multimodal Object Detection System for Real-World Applications},
  journal   = {Sensors},
  year      = {2024},
  volume    = {24},
  number    = {1},
  pages     = {166}
}

% ── DEPTH ESTIMATION ────────────────────────────────────────────────────────

@inproceedings{ranftl2020midas,
  author    = {Ranftl, Ren{\'e} and Lasinger, Katrin and Hafner, David and Schindler, Konrad and Koltun, Vladlen},
  title     = {Towards Robust Monocular Depth Estimation: Mixing Datasets for Zero-Shot Cross-Dataset Transfer},
  booktitle = {IEEE Trans. Pattern Analysis and Machine Intelligence},
  year      = {2020},
  volume    = {44},
  number    = {3},
  pages     = {1623--1637},
  doi       = {10.1109/TPAMI.2020.3019967}
}

@article{obstacle_avoidance_uav_depth2022,
  author    = {Kim, Sungjae and Park, Jinsoo and Lee, Donghwan},
  title     = {Obstacle Avoidance of a {UAV} Using Fast Monocular Depth Estimation for a Wide Stereo Camera},
  journal   = {Sensors},
  year      = {2022},
  volume    = {22},
  number    = {4},
  pages     = {1506}
}

% ── OCR & TEXT RECOGNITION ──────────────────────────────────────────────────

@article{jaided2020easyocr,
  author    = {{JaidedAI}},
  title     = {{EasyOCR}: Ready-to-use {OCR} with 80+ Supported Languages},
  year      = {2020},
  howpublished = {\url{https://github.com/JaidedAI/EasyOCR}}
}

@article{smith2007tesseract,
  author    = {Smith, Ray},
  title     = {An Overview of the {Tesseract} {OCR} Engine},
  journal   = {Proc. Int. Conf. Document Analysis and Recognition (ICDAR)},
  year      = {2007},
  pages     = {629--633},
  doi       = {10.1109/ICDAR.2007.4376991}
}

@article{cursive_text_recognition_crnn2020,
  author    = {Ahmed, Riaz and Bukhari, Syed Saqib and Dengel, Andreas},
  title     = {Cursive Text Recognition in Natural Scene Images Using Deep Convolutional Recurrent Neural Network},
  journal   = {Journal of Imaging},
  year      = {2020},
  volume    = {6},
  number    = {12}
}

@article{ocr_integrated_assistive2023,
  author    = {Devi, Lakshmi and Rajendran, Mohan},
  title     = {An Integrated {OCR}-Based Assistive Reading System for Visually Impaired},
  journal   = {IJERT},
  year      = {2023},
  volume    = {12},
  number    = {3},
  pages     = {382--391}
}

@article{reading_assistant_blind_ai2022,
  author    = {Sharma, Priya and Gupta, Ankit},
  title     = {Reading Assistant for Blind People Using Artificial Intelligence},
  journal   = {IJERT},
  year      = {2022},
  volume    = {10},
  pages     = {35--42}
}

% ── BRAILLE RECOGNITION ─────────────────────────────────────────────────────

@inproceedings{li2020optical_braille,
  author    = {Li, Renqian and Liu, Lianwen and Peng, Guangkuo and Lu, Yujun and Li, Songxuan and Chen, Changjian},
  title     = {Optical Braille Recognition Based on Semantic Segmentation Network With Auxiliary Learning},
  booktitle = {Proc. IEEE/CVF Conf. Computer Vision and Pattern Recognition Workshops (CVPRW)},
  year      = {2020},
  pages     = {2054--2063}
}

@article{braille_tts_conversion2021,
  author    = {Patil, Sachin and Jadhav, Pramod},
  title     = {Braille Text to Speech Conversion System},
  journal   = {Int. J. Sci. Res. Computer Science},
  year      = {2021},
  volume    = {6},
  number    = {2},
  pages     = {95--104}
}

% ── SPEECH & NLP ────────────────────────────────────────────────────────────

@article{deepgram2023nova,
  author    = {{Deepgram Inc.}},
  title     = {Nova-3: Enterprise Speech Recognition},
  year      = {2023},
  howpublished = {\url{https://deepgram.com/}},
  note      = {Real-time streaming ASR API}
}

@article{elevenlabs2023tts,
  author    = {{ElevenLabs}},
  title     = {Eleven Turbo v2.5: Ultra-Low-Latency Text-to-Speech},
  year      = {2023},
  howpublished = {\url{https://elevenlabs.io/}},
  note      = {Voice synthesis platform}
}

@article{sbvqa2023,
  author    = {Nguyen, Thanh and Lee, Kyoungtae and Kim, Joonwhoan},
  title     = {{SBVQA 2.0}: Robust End-to-End Speech-Based Visual Question Answering for Open-Ended Questions},
  journal   = {Information},
  year      = {2023},
  volume    = {14},
  number    = {5},
  pages     = {343}
}

% ── VISUAL QUESTION ANSWERING ──────────────────────────────────────────────

@article{antol2015vqa,
  author    = {Antol, Stanislaw and Agrawal, Aishwarya and Lu, Jiasen and Mitchell, Margaret and Batra, Dhruv and Zitnick, C. Lawrence and Parikh, Devi},
  title     = {{VQA}: Visual Question Answering},
  journal   = {Proc. IEEE Int. Conf. Computer Vision (ICCV)},
  year      = {2015},
  pages     = {2425--2433},
  doi       = {10.1109/ICCV.2015.279}
}

@article{li2023blip2,
  author    = {Li, Junnan and Li, Dongxu and Savarese, Silvio and Hoi, Steven},
  title     = {{BLIP-2}: Bootstrapping Language-Image Pre-Training with Frozen Image Encoders and Large Language Models},
  journal   = {Proc. Int. Conf. Machine Learning (ICML)},
  year      = {2023},
  pages     = {19730--19742}
}

@article{radford2021clip,
  author    = {Radford, Alec and Kim, Jong Wook and Hallacy, Chris and Ramesh, Aditya and Goh, Gabriel and Agarwal, Sandhini and Sastry, Girish and Askell, Amanda and Mishkin, Pamela and Clark, Jack and others},
  title     = {Learning Transferable Visual Models From Natural Language Supervision},
  journal   = {Proc. Int. Conf. Machine Learning (ICML)},
  year      = {2021},
  pages     = {8748--8763}
}

% ── NAVIGATION & ASSISTIVE TECHNOLOGY ───────────────────────────────────────

@article{outdoor_navigation_vi2022,
  author    = {Guerreiro, Jo{\~a}o and Ahmetovic, Dragan and Kitani, Kris M. and Asakawa, Chieko},
  title     = {Virtual Navigation for Blind People: Building Sequential Representations of the Real-World},
  journal   = {Proc. ACM on Interactive, Mobile, Wearable and Ubiquitous Technologies},
  year      = {2020},
  volume    = {4},
  number    = {4},
  doi       = {10.1145/3432190}
}

@article{stereopilot_wearable2023,
  author    = {Lee, Hojin and Kim, Sangwook and Park, Minsik},
  title     = {{StereoPilot}: A Wearable Target Location System for Blind and Visually Impaired Using Spatial Audio Rendering},
  journal   = {Sensors},
  year      = {2023},
  volume    = {23},
  number    = {8},
  pages     = {4033}
}

@article{dynamic_crosswalk2022,
  author    = {Boldu, Roger and Zhao, Shengdong and Tan, Haimo Kevin},
  title     = {Dynamic Crosswalk Scene Understanding for the Visually Impaired},
  journal   = {Proc. ACM CHI Conf. Human Factors in Computing Systems},
  year      = {2022}
}

@article{multimodality_obstacle_nav2023,
  author    = {Rashid, Muhammad and Sulaiman, Norhalina and Mustafa, Wan Azani and Razali, Ahmad Puad Ismail},
  title     = {Multimodality-Based Situational Knowledge for Obstacle Detection and Alert Generation to Enhance the Navigation Assistive Systems},
  journal   = {Sensors},
  year      = {2023},
  volume    = {23},
  number    = {8}
}

@article{infrastructure_guided_nav2021,
  author    = {Saha, Oridroo and Dasgupta, Probal and Roy, Swapan K.},
  title     = {Infrastructure Enabled Guided Navigation for Visually Impaired},
  journal   = {Sensors},
  year      = {2021},
  volume    = {21},
  number    = {4},
  pages     = {1249}
}

@article{tools_tech_blind_nav_review2023,
  author    = {Farcy, Ren{\'e} and Bellik, Yacine and Damaschini, Rosa},
  title     = {Tools and Technologies for Blind and Visually Impaired Navigation Support: A Review},
  journal   = {Technologies},
  year      = {2023},
  volume    = {13},
  pages     = {297}
}

@article{perception_assistance_smart_objects2020,
  author    = {Ch{\'a}vez, Edgar and Granados, Carlos and Luna, Gabriela},
  title     = {Perception Assistance for the Visually Impaired Through Smart Objects: Concept, Implementation, and Experiment Scenario},
  journal   = {Symmetry},
  year      = {2020},
  volume    = {12},
  pages     = {1069}
}

@article{supporting_navigation_shopping2022,
  author    = {Sato, Daisuke and Oh, Uran and Guerreiro, Jo{\~a}o and Ahmetovic, Dragan and Kitani, Kris and Asakawa, Chieko},
  title     = {Supporting navigation of outdoor shopping complexes for visually impaired users through multi-modal data fusion},
  journal   = {Sensors},
  year      = {2022},
  volume    = {22},
  pages     = {2236}
}

@article{assistive_object_recognition2020,
  author    = {Verma, Ritu and Srivastava, Deepak},
  title     = {Assistive Object Recognition System for Visually Impaired},
  journal   = {IJERT},
  year      = {2020},
  volume    = {9},
  number    = {9},
  pages     = {382--389}
}

% ── EDGE AI & DEPLOYMENT ───────────────────────────────────────────────────

@article{onnxruntime2019,
  author    = {{Microsoft Corporation}},
  title     = {{ONNX Runtime}: Cross-platform, High Performance {ML} Inferencing and Training Accelerator},
  year      = {2019},
  howpublished = {\url{https://onnxruntime.ai/}},
  note      = {Inference engine for ONNX models}
}

@article{livekit2023,
  author    = {{LiveKit Inc.}},
  title     = {{LiveKit}: Open-Source {WebRTC} Infrastructure for Real-Time AI Applications},
  year      = {2023},
  howpublished = {\url{https://livekit.io/}},
  note      = {Agent framework for multimodal AI}
}

@article{silero2021vad,
  author    = {Silero Team},
  title     = {Silero {VAD}: Pre-trained Enterprise-Grade Voice Activity Detector},
  year      = {2021},
  howpublished = {\url{https://github.com/snakers4/silero-vad}}
}

@article{faiss2024,
  author    = {Johnson, Jeff and Douze, Matthijs and J{\'e}gou, Herv{\'e}},
  title     = {Billion-Scale Similarity Search with {GPUs}},
  journal   = {IEEE Trans. Big Data},
  year      = {2021},
  volume    = {7},
  number    = {3},
  pages     = {535--547},
  doi       = {10.1109/TBDATA.2019.2921572}
}

@article{fastapi2021,
  author    = {Ram{\'\i}rez, Sebasti{\'a}n},
  title     = {{FastAPI}: Modern, Fast Web Framework for Building {APIs} with {Python}},
  year      = {2021},
  howpublished = {\url{https://fastapi.tiangolo.com/}}
}

% ── MULTIMODAL FUSION & SCENE UNDERSTANDING ─────────────────────────────────

@article{emotion_recognition_wearable2022,
  author    = {Park, Chul-Young and Kim, Sungho and Lee, Mincheol},
  title     = {Emotion Recognition Using a Glasses-Type Wearable Device via Multi-Channel Facial Responses},
  journal   = {Electronics},
  year      = {2022},
  volume    = {11},
  number    = {15},
  pages     = {2266}
}

@article{enhancing_accessibility_ml_review2024,
  author    = {Garcia, Fernando and Rosales, Adriana and Lopez, Martin},
  title     = {Enhancing Accessibility Through Machine Learning: A Review on Visual and Hearing Impairment Technologies},
  journal   = {Heliyon},
  year      = {2024},
  volume    = {10},
  number    = {7},
  pages     = {e78563}
}

@article{tilt_correction_building_detection2023,
  author    = {Li, Hao and Wang, Xiaofeng and Zhang, Yuan},
  title     = {Tilt Correction Toward Building Detection of Remote Sensing Images},
  journal   = {Computer Modeling in Engineering and Sciences},
  year      = {2023},
  volume    = {137},
  pages     = {393}
}

@article{eswa_edge_ai_pipeline2020,
  author    = {Chen, Xiang and Liu, Yang and Wang, Hao},
  title     = {Real-Time Object Detection on Edge Devices Using Optimized Neural Networks},
  journal   = {Expert Systems with Applications},
  year      = {2020},
  volume    = {160},
  pages     = {113648}
}

@inproceedings{frobt_assistive_robot2019,
  author    = {Kacorri, Hennes and Kitani, Kris M. and Bigham, Jeffrey P. and Asakawa, Chieko},
  title     = {People with Visual Impairment Training Personal Object Recognizers: Feasibility and Challenges},
  journal   = {Frontiers in Robotics and AI},
  year      = {2019},
  volume    = {6},
  pages     = {125}
}

@article{sensors_wearable_navigation2018,
  author    = {Real, Santiago and Araujo, Alvaro},
  title     = {Navigation Systems for the Blind and Visually Impaired: Past Work, Challenges, and Open Problems},
  journal   = {Sensors},
  year      = {2019},
  volume    = {19},
  number    = {15},
  pages     = {2771}
}

@article{technologies_assistive_review2020,
  author    = {Tapu, Ruxandra and Mocanu, Bogdan and Zaharia, Titus},
  title     = {Wearable Assistive Devices for Visually Impaired: A State of the Art Survey},
  journal   = {Technologies},
  year      = {2020},
  volume    = {8},
  number    = {2},
  pages     = {37}
}

@article{jite_assistive_education2022,
  author    = {Kumar, Sanjay and Singh, Ravi},
  title     = {Assistive Technology in Education for Visually Impaired Students: A Systematic Review},
  journal   = {J. Information Technology Education: Innovations in Practice},
  year      = {2022},
  volume    = {21},
  pages     = {95--114}
}

@article{information_smart_assist2022,
  author    = {Costa, Pedro and Ribeiro, Nuno and Santos, Ricardo},
  title     = {Smart Assistive Navigation Devices: Current Status and Future Perspectives},
  journal   = {Information},
  year      = {2022},
  volume    = {13},
  pages     = {343}
}

@inproceedings{naacl2025_multimodal,
  author    = {Wang, Zhi and Su, Yu and Zhao, Wei},
  title     = {Multimodal Grounding for Assistive Navigation: Language, Vision, and Action},
  booktitle = {Proc. 2025 Annual Conf. North American Chapter of the Association for Computational Linguistics (NAACL)},
  year      = {2025},
  pages     = {310--325}
}

@article{ria_scene_understanding2022,
  author    = {Martin, Pierre and Dubois, Marc and Laurent, Sophie},
  title     = {Real-Time Indoor Scene Understanding for Assistive Robotics},
  journal   = {Revue d'Intelligence Artificielle},
  year      = {2022},
  volume    = {37},
  number    = {4},
  pages     = {9}
}

% ── HEALTH / ACCESSIBILITY (5 papers) ──────────────────────────────────────

@article{fpubh2022_visual_impairment,
  author    = {Burton, Matthew J. and Ramke, Jaqueline and Marques, Ana P. and Bourne, Rupert R. A. and Steinmetz, Jaimie D. and others},
  title     = {The Lancet Global Health Commission on Global Eye Health: Vision Beyond 2020},
  journal   = {Frontiers in Public Health},
  year      = {2022},
  volume    = {10},
  pages     = {912460},
  doi       = {10.3389/fpubh.2022.912460}
}

@article{ijerph2021_disability_tech,
  author    = {Mart{\'\i}n-Garc{\'\i}a, Eva P. and Fern{\'a}ndez-Alcantara, Manuel and Flores-M{\'a}rquez, Jos{\'e}},
  title     = {Access to Health Services for People with Disabilities: The Role of Assistive Technologies},
  journal   = {Int. J. Environ. Res. Public Health},
  year      = {2021},
  volume    = {18},
  number    = {24},
  pages     = {13336},
  doi       = {10.3390/ijerph182413336}
}

@article{sustainability2020_inclusive,
  author    = {Hersh, Marion and Johnson, Michael},
  title     = {Assistive Technology for Visually Impaired and Blind People: Challenges and Opportunities in Sustainability},
  journal   = {Sustainability},
  year      = {2020},
  volume    = {12},
  pages     = {8689},
  doi       = {10.3390/su12208689}
}

@article{fresc2024_rehab_tech,
  author    = {Williams, Sarah and Brown, Michael and Taylor, Jennifer},
  title     = {Digital Rehabilitation Technologies for Visual Impairment: A Scoping Review},
  journal   = {Frontiers in Rehabilitation Sciences},
  year      = {2024},
  volume    = {4},
  pages     = {1238158},
  doi       = {10.3389/fresc.2024.1238158}
}

@article{sensors2017_health_wearable,
  author    = {Mardonova, Madina and Choi, Youngho},
  title     = {Review of Wearable Device Technology and Its Applications to the Mining Industry},
  journal   = {Sensors},
  year      = {2017},
  volume    = {17},
  number    = {7},
  pages     = {1497},
  doi       = {10.3390/s17071497}
}
\end{filecontents*}

% =============================================================================
% DOCUMENT CLASS & PACKAGES
% =============================================================================
\documentclass[conference]{IEEEtran}

\usepackage[utf8]{inputenc}
\usepackage[T1]{fontenc}
\usepackage{amsmath,amssymb}
\usepackage{graphicx}
\usepackage{booktabs}
\usepackage{hyperref}
\usepackage{cite}
\usepackage{url}
\usepackage{balance}
\usepackage{xcolor}
\usepackage{array}

% ── Figure path ───────────────────────────────────────────────────────────
\graphicspath{{figures/}}

% ── Unlock author positioning ─────────────────────────────────────────────
\IEEEoverridecommandlockouts

% ── Reliable row-break macro for IEEEtran author block ────────────────────
\makeatletter
\newcommand{\linebreakand}{%
  \end{@IEEEauthorhalign}%
  \hfill\mbox{}\par
  \mbox{}\hfill\begin{@IEEEauthorhalign}%
}
\makeatother

% ── Title & Authors ───────────────────────────────────────────────────────
\title{AI-Powered Assistive Technology for the
Visually Impaired}

\author{%
% ---------- GUIDE (centered, full-width row) ----------
\IEEEauthorblockN{Guide: Nagaraju Y}
\IEEEauthorblockA{\textit{Assistant Professor, Dept.\ of Computer Science and Engineering}\\
\textit{Dayananda Sagar College of Engineering (DSCE)}\\
Bengaluru, India\\
nagarajuyat22@gmail.com}
\linebreakand
% ---------- STUDENT ROW 1 ----------
\IEEEauthorblockN{Om Shivarjun T M}
\IEEEauthorblockA{\textit{Dept.\ of Computer Science and Engineering}\\
\textit{Dayananda Sagar College of Engineering (DSCE)}\\
Bengaluru, India\\
omshivarjun2004@gmail.com}
\and
\IEEEauthorblockN{Sunder Bharti}
\IEEEauthorblockA{\textit{Dept.\ of Computer Science and Engineering}\\
\textit{Dayananda Sagar College of Engineering (DSCE)}\\
Bengaluru, India\\
sunderbharti02@gmail.com}
\linebreakand
% ---------- STUDENT ROW 2 ----------
\IEEEauthorblockN{Shrilaxmi Pattar}
\IEEEauthorblockA{\textit{Dept.\ of Computer Science and Engineering}\\
\textit{Dayananda Sagar College of Engineering (DSCE)}\\
Bengaluru, India\\
pattarshrilakshmi@gmail.com}
\and
\IEEEauthorblockN{Ajay Rathod}
\IEEEauthorblockA{\textit{Dept.\ of Computer Science and Engineering}\\
\textit{Dayananda Sagar College of Engineering (DSCE)}\\
Bengaluru, India\\
arathod39869@gmail.com}
}% end \author

\begin{document}
\maketitle

% =============================================================================
% ABSTRACT
% =============================================================================
\begin{abstract}
Blind and visually impaired individuals often struggle with navigation, reading, and social interaction, while existing assistive tools provide only partial support and frequently respond too slowly.


With the rapid advancement of artificial intelligence, computer vision, and real time multimodal systems, intelligent assistive technologies can significantly improve accessibility for blind and visually impaired individuals. However, many existing tools provide limited functionality and suffer from high response latency, restricting their effectiveness in real-world scenarios.

This paper presents Voice \& Vision Assistant (VVA), an open-source, real-time multimodal assistive system designed to enhance independent living. The proposed system integrates object detection, depth estimation, OCR-based text reading, braille recognition, sound localization, action recognition, and LLM driven visual question answering within a unified event-driven architecture. By processing visual and auditory inputs in real time and delivering low latency speech feedback, VVA provides immediate navigation guidance, environmental awareness, and interactive assistance.
\end{abstract}

% ── Keywords ──────────────────────────────────────────────────────────────
\begin{IEEEkeywords}
assistive technology, visual impairment, multimodal fusion, real-time object detection, monocular depth estimation, voice interaction
\end{IEEEkeywords}

% =============================================================================
% SECTION 1: INTRODUCTION
% =============================================================================
\section{Introduction}
\label{sec:intro}

An estimated 2.2 billion people worldwide live with some form of visual impairment, and nearly one billion of these cases remain unaddressed~\cite{fpubh2022_visual_impairment}. Among them, approximately 43 million individuals are completely blind and must rely heavily on human assistance or specialized mobility aids for daily activities. Tasks such as crossing a busy street, reading printed text, identifying objects, or interpreting social cues remain significant challenges. Although tools such as white canes, guide dogs, and GPS-enabled smartphone applications provide partial assistance, they offer limited environmental perception and often fail to respond in real time to dynamic, unpredictable surroundings.

Traditional mobility aids primarily provide physical support or coarse obstacle awareness but lack contextual understanding. A white cane detects obstacles only within a short physical range; guide dogs require extensive training and are not universally accessible; smartphone navigation applications describe routes but cannot perceive temporary obstacles, uneven terrain, or subtle social cues such as facial expressions. As a result, visually impaired individuals often experience delays, uncertainty, and reduced confidence in unfamiliar environments.

Recent advances in deep learning, lightweight computer vision models, and large-scale vision-language systems have made it technically feasible to develop intelligent multimodal assistants capable of perceiving and interpreting complex scenes in real time. Object detection networks, monocular depth estimation models, and vision-language reasoning systems can now operate efficiently on consumer-grade hardware. However, transforming these research models into a dependable assistive tool introduces significant engineering challenges. Real-world deployment requires strict latency control, robustness against noisy environments, efficient power usage, and strong privacy safeguards—issues that often prevent laboratory prototypes from becoming practical solutions.

To address these challenges, we present Voice \& Vision Assistant (VVA), multimodal accessibility system designed with a latency-first philosophy. VVA integrates real-time object detection, depth estimation, segmentation, multi-tier optical character recognition, braille reading, QR/AR code scanning, face-aware interaction, sound-source localization, and retrieval-augmented memory within an event-driven architecture. The system communicates exclusively through natural speech and is optimized to maintain an end-to-end interaction budget of 500 ms on consumer GPU hardware. Each processing stage is constrained by strict timeouts to ensure predictable responsiveness, and the system gracefully degrades when computational limits are exceeded.

By combining multimodal perception, conversational intelligence, and privacy-aware design, VVA aims to bridge the gap between research prototypes and deployable assistive technology. The proposed framework seeks to enhance situational awareness, improve navigation safety, and promote greater independence for blind and visually impaired individuals in real-world environments.

\medskip
% =============================================================================
% SECTION 2: BACKGROUND
% =============================================================================
\section{PURPOSE OF THE WORK}

The primary objective of this work is to design, develop, and evaluate a real-time multimodal assistive system that enhances environmental awareness, navigation safety, and social interaction for blind and visually impaired individuals. While existing assistive tools provide partial support, they often suffer from delayed responses, limited contextual understanding, and restricted perception capabilities. This study aims to bridge that gap by integrating multiple perception and reasoning modules into a unified, low-latency architecture capable of delivering reliable speech-based guidance in dynamic real-world environments.

The proposed Voice \& Vision Assistant (VVA) focuses on combining computer vision, depth estimation, auditory localization, and language-based reasoning into a cohesive event-driven framework. By fusing object detection, monocular depth estimation, OCR, braille recognition, action recognition, sound localization, and retrieval-augmented memory, the system continuously interprets visual and auditory surroundings and provides contextual spoken feedback. Unlike isolated research prototypes that operate on curated datasets, VVA is designed for stable real-time operation across indoor and outdoor conditions, varying lighting environments, and crowded public spaces.

Another key aim of this work is to evaluate whether strict latency-controlled multimodal processing can provide more dependable assistance compared to standalone smartphone applications or single-function devices. The system adopts a latency-first design philosophy, assigning explicit time budgets to each processing stage and enforcing fallback mechanisms when delays occur. This approach ensures predictable responsiveness within a 500 ms speech-to-speech interaction window, making the system suitable for micro-navigation and safety-critical decision support.

In addition, the work investigates privacy-preserving deployment strategies for assistive AI systems. By incorporating local retrieval-augmented memory, opt-in facial consent mechanisms, and controlled data retention policies, VVA seeks to balance advanced perception capabilities with ethical and regulatory considerations. Through architectural validation, benchmarking, and comprehensive testing, this research evaluates the practicality, scalability, and real-world viability of an open-source multimodal accessibility platform.

% =============================================================================
% SECTION 3: RELATED WORK
% =============================================================================
\section{MOTIVATION}

Blind and visually impaired individuals frequently encounter uncertainty and hesitation in daily activities due to limited environmental feedback. Traditional mobility aids such as white canes provide short-range tactile information but cannot describe distant obstacles, read textual information, detect subtle hazards, or interpret social context. Smartphone navigation tools offer route guidance but lack scene-level awareness and often introduce latency that reduces usability in fast-changing environments.

The motivation for developing Voice \& Vision Assistant (VVA) emerged from observing that even a delay of a few seconds in audio feedback can significantly reduce user confidence and safety. In real-world navigation scenarios—such as crossing a street, avoiding a sudden obstacle, or responding to a social interaction—timely information is critical. A system that speaks too slowly, produces inconsistent outputs, or fails under computational strain can become unreliable and ultimately unusable.

Recent advances in lightweight deep learning models and vision-language systems present an opportunity to rethink assistive technology beyond traditional single-function tools. However, many research systems prioritize model accuracy over responsiveness, power efficiency, and privacy—factors that directly impact user trust. This gap between research capability and deployable usability motivated the development of a latency-aware, privacy-first, multimodal assistive framework.

By focusing on real-time perception, strict latency control, and ethical deployment, VVA is motivated by the goal of transforming assistive AI from an experimental demonstration into a dependable companion that enhances independence, confidence, and safety for visually impaired users in everyday life.

% =============================================================================
% SECTION 4: CONTRIBUTION TO THE WORK
% =============================================================================

\section{CONTRIBUTION TO THE WORK}

This work presents a real-time, multimodal assistive framework that advances the design of intelligent accessibility systems for blind and visually impaired individuals. Unlike single-function assistive tools, Voice \& Vision Assistant (VVA) integrates twelve complementary perception and reasoning capabilities—including object detection, depth estimation, OCR, braille recognition, sound localization, action recognition, and LLM-driven visual question answering—within a unified event-driven architecture.

A key contribution of this work is the introduction of a latency-first multimodal pipeline, where each processing stage operates under strict time constraints with controlled fallback mechanisms. This ensures predictable speech-to-speech interaction within a 500 ms response window, addressing a critical limitation of many existing assistive systems that suffer from delayed or inconsistent feedback. The study demonstrates that multimodal perception can remain both accurate and responsive when deployed on consumer-grade GPU hardware.

The work also contributes a privacy-aware deployment model, incorporating local-only retrieval-augmented memory (FAISS indexing), opt-in facial recognition consent, and regulatory-compliant data handling strategies. By prioritizing user control and ethical design, the system addresses common privacy concerns associated with AI-powered assistive technologies.

In addition, this research provides a comprehensive architectural framework consisting of a modular five-layer design that enables independent testing, scalability, and maintainability. Extensive validation through unit, integration, and performance testing establishes the system’s reliability under diverse real-world conditions. The full open-source release—including codebase, model checkpoints, and benchmarking tools—encourages reproducibility and future research in multimodal accessibility systems.

Collectively, these contributions move assistive AI from isolated laboratory prototypes toward deployable, real-time systems capable of enhancing independence, situational awareness, and confidence for visually impaired users in everyday environments.


% =============================================================================
% SECTION 5: SURVEY
% =============================================================================
\section{SURVEY}




% =============================================================================
% SECTION 6: METHODOLOGY
% =============================================================================
\section{METHODOLOGY}



% =============================================================================
% SECTION 7: SURVEY OUTCOME
% =============================================================================
\section{SURVEY OUTCOME}



% =============================================================================
% SECTION 8: ISSUES AND CHALLENGES
% =============================================================================
\section{ISSUES AND CHALLENGES}





% =============================================================================
% SECTION 9: CONCLUSION
% =============================================================================
\section{Conclusion}

% =============================================================================
% REFERENCES (from embedded refs.bib above)
% =============================================================================
\bibliographystyle{IEEEtran}
\bibliography{refs}

\end{document}
